\scanpage{342}
% top centre
\mspage{258}
% lines 1–9
est de cardinal $\leq\pi$, et que la catégorie $\unsure{\mathcal E|\mathscr E|\mathcal C}_j$ est \unsure{grde|gde} devant $\pi$, on voit qu'on peut bel et bien trouver $X_{\alpha+\unsure{1}}(j)$ et des $t_\alpha(\unsure{\lambda|j}):X_\alpha(j')\unsure{\longrightarrow|\mathrel{-}}X_{\alpha+1}(j)$, de façon à satisfaire (si $\alpha=\unsure{\alpha'|\beta'}+2$) aux contraintes \unsure{explicitées} du $2^\circ$. Cela achève la \deleted{\unsure{On doit}}\inserted{description de} la construction de $X_{\alpha+1}(\deleted{\unsure{j}}\unsure{\cdot|j})$, en termes des \deleted{contraintes} $X_\beta(\cdot)$ $(\beta\leq\alpha)$ et des $t_\beta(\lambda)$ $(\beta<\alpha)$.

% lines 10–17
Il faut enfin définir $X_\alpha(\cdot)$ pour $\alpha$ un ordinal limite. Cette fois, on lui demande seulement de recevoir les $X_\beta(\deleted{\unsure{j}}\cdot)$ pour $\beta<\alpha$. Comme l'ens $\unsure{I|J}_\alpha$ de ces $\beta$ est de cardinal $\leq\pi$ et que les $\mathcal E_j$ sont \unsure{grdes|gdes} devant $\pi$, on \unsure{p}\deleted{\unsure{eut}} bel et bien trouver les $X_\alpha(j)$\unsure{.|:}
% archive sheet number: 341
